% experiments_reasonspec.tex — Experiments section for Reasoning Speculation
% Included from main_reasonspec.tex via % experiments_reasonspec.tex — Experiments section for Reasoning Speculation
% Included from main_reasonspec.tex via % experiments_reasonspec.tex — Experiments section for Reasoning Speculation
% Included from main_reasonspec.tex via % experiments_reasonspec.tex — Experiments section for Reasoning Speculation
% Included from main_reasonspec.tex via \input{sections/experiments_reasonspec}

% ---------------------------------------------------------------------------
\subsection{Experimental Setup}
\label{sec:exp:setup}
% ---------------------------------------------------------------------------

\paragraph{Model and inference.}
All experiments use \textbf{Qwen3-8B} with thinking mode enabled
(\texttt{enable\_thinking=True}), which instructs the model to produce an
extended chain-of-thought within \texttt{<think>...</think>} tags before
emitting a final answer.
We control the thinking token budget $b$ by setting
\texttt{max\_thinking\_tokens}; once this limit is reached, the model is
forced to exit the thinking block and produce an answer.
Greedy decoding ($\tau = 0$) is used for the first exploration path and
all resolution phases; the remaining $K{-}1$ exploration paths use
$\tau = 0.7$ to encourage diversity (\S\ref{sec:method:explore}).

\paragraph{Benchmark.}
We evaluate on \textbf{GSM8K}~\citep{cobbe2021training}, a dataset of
1{,}319 grade-school math word problems.
For the main experiments and ablations we sample subsets of $n = 200$
problems (with fixed random seeds for reproducibility); full-dataset
results are reported for the best configuration.
Accuracy is measured via exact match after deterministic number extraction
(the rightmost numerical value in the model's output).

\paragraph{Baselines.}
We compare \method{} against four baselines that span the natural
design space:
\begin{itemize}[nosep,leftmargin=*]
  \item \textbf{Fixed-$b$}: single-pass generation with a fixed thinking
    budget $b \in \{128, 256, 512\}$ tokens.
  \item \textbf{SC@$K{\times}128$}: self-consistency~\citep{wang2023selfconsistency}
    with $K$ independent samples at $b = 128$ tokens each, followed by
    majority voting.
    This baseline uses the same total exploration tokens as \method{}
    with $K$ paths and $B_{\mathrm{probe}} = 128$, providing a
    compute-matched comparison.
\end{itemize}

\paragraph{\method{} configurations.}
We explore the following hyperparameter axes:
parallelism $K \in \{2, 3, 4, 8\}$,
probe budget $B_{\mathrm{probe}} \in \{128, 170, 256, 512\}$,
medium resolution budget $B_{\mathrm{med}} \in \{256, 512\}$, and
hard resolution budget $B_{\mathrm{hard}} \in \{512, 1024\}$.
Routing thresholds default to $\theta_E = 0.75$ (easy) and
$\theta_M = 0.5$ (medium) unless varied in ablation studies.
The consensus confidence weights are fixed at
$\gamma = 0.5\alpha + 0.3\nu + 0.2\mu$ throughout.

\paragraph{Evaluation metrics.}
For each method we report:
(1)~\textbf{Accuracy} (exact-match \%),
(2)~\textbf{Avg Tokens} (mean thinking tokens consumed per sample,
including exploration overhead for \method{}),
(3)~\textbf{Route distribution} (fraction of samples classified as
Easy / Medium / Hard by the consensus router).

\paragraph{Hardware.}
All experiments are conducted on a single NVIDIA A100-80GB GPU.

% ---------------------------------------------------------------------------
\subsection{Main Results}
\label{sec:exp:main}
% ---------------------------------------------------------------------------

Table~\ref{tab:main-rs} presents the primary comparison between
\method{}, fixed-budget baselines, and self-consistency on GSM8K
($n = 200$).
All \method{} variants use default routing thresholds
($\theta_E = 0.75$, $\theta_M = 0.5$).

\begin{table}[t]
\centering
\caption{%
  \textbf{Main results on GSM8K} ($n = 200$, Qwen3-8B thinking mode).
  \textbf{Avg Tok} includes full end-to-end cost (exploration +
  resolution).
  Route columns report the fraction of samples classified as Easy (E),
  Medium (M), and Hard (H) by the consensus router.
  Best adaptive result in \textbf{bold}; best baseline \underline{underlined}.
}
\label{tab:main-rs}
\small
\setlength{\tabcolsep}{4pt}
\begin{tabular}{@{}lcccccccc@{}}
\toprule
\textbf{Method} & $K$ & $B_{\mathrm{probe}}$ &
  \textbf{Acc (\%)} & \textbf{Avg Tok} &
  \textbf{E\%} & \textbf{M\%} & \textbf{H\%} \\
\midrule
\multicolumn{8}{@{}l}{\textit{Fixed-budget baselines}} \\
Fixed-128       & -- & -- & XX.X & 128  & -- & -- & -- \\
Fixed-256       & -- & -- & XX.X & 256  & -- & -- & -- \\
\underline{Fixed-512} & -- & -- & \underline{XX.X} & 512  & -- & -- & -- \\
\midrule
\multicolumn{8}{@{}l}{\textit{Self-consistency baselines}} \\
SC@$2{\times}128$  & 2 & 128 & XX.X & 256  & -- & -- & -- \\
SC@$4{\times}128$  & 4 & 128 & XX.X & 512  & -- & -- & -- \\
SC@$8{\times}128$  & 8 & 128 & XX.X & 1024 & -- & -- & -- \\
\midrule
\multicolumn{8}{@{}l}{\textit{\method{} (ours)}} \\
\method{}-2      & 2 & 128 & XX.X & XX.X & XX.X & XX.X & XX.X \\
\method{}-4      & 4 & 128 & XX.X & XX.X & XX.X & XX.X & XX.X \\
\method{}-4-256  & 4 & 256 & XX.X & XX.X & XX.X & XX.X & XX.X \\
\textbf{\method{}-8} & 8 & 128 & \textbf{XX.X} & XX.X & XX.X & XX.X & XX.X \\
\method{}-8-256  & 8 & 256 & XX.X & XX.X & XX.X & XX.X & XX.X \\
\bottomrule
\end{tabular}
\end{table}

Several observations emerge from Table~\ref{tab:main-rs}.
First, \method{} consistently outperforms fixed-budget baselines at
comparable or lower token cost.
The best configuration (\method{}-8) achieves XX.X\% accuracy using an
average of XX.X tokens---a \textbf{+XX.X\,pp} improvement over Fixed-512
at \textbf{XX\%} of the token budget.
Second, \method{} substantially outperforms the compute-matched
self-consistency baseline (SC@$K{\times}128$), confirming that the
consensus signal is more effectively leveraged through adaptive routing
than through simple majority voting.
For instance, at $K = 4$, \method{}-4 achieves XX.X\% vs.\
SC@$4{\times}128$'s XX.X\%, while consuming fewer total tokens because
easy samples are resolved without a resolution phase.
Third, the route distribution reveals that a substantial fraction of
GSM8K problems (XX.X\%) are routed as easy, validating the central
premise that many problems can be solved with minimal compute.

% ---------------------------------------------------------------------------
\subsection{Ablation Studies}
\label{sec:exp:ablation}
% ---------------------------------------------------------------------------

We systematically ablate the three key hyperparameters of \method{} to
understand their individual contributions.

% --- K ablation ---
\paragraph{Effect of parallelism $K$.}
Table~\ref{tab:ablation-k} fixes $B_{\mathrm{probe}} = 128$ and varies
$K \in \{2, 3, 4, 8\}$.
Increasing $K$ improves the reliability of the consensus signal
(Proposition~\ref{prop:monotone}), yielding monotonically higher accuracy.
However, the exploration cost scales linearly in $K$, creating a
diminishing-returns trade-off: from $K = 4$ to $K = 8$, accuracy improves
by XX.X\,pp while average tokens increase by XX.X.
We find $K = 4$ to be the best cost-accuracy operating point for GSM8K.

\begin{table}[t]
\centering
\caption{%
  \textbf{Ablation on $K$} (parallelism) with $B_{\mathrm{probe}} = 128$,
  $B_{\mathrm{med}} = 256$, $B_{\mathrm{hard}} = 512$, default thresholds.
}
\label{tab:ablation-k}
\small
\begin{tabular}{@{}ccccccc@{}}
\toprule
$K$ & \textbf{Acc (\%)} & \textbf{Avg Tok} &
  \textbf{E\%} & \textbf{M\%} & \textbf{H\%} \\
\midrule
2 & XX.X & XX.X & XX.X & XX.X & XX.X \\
3 & XX.X & XX.X & XX.X & XX.X & XX.X \\
4 & XX.X & XX.X & XX.X & XX.X & XX.X \\
8 & XX.X & XX.X & XX.X & XX.X & XX.X \\
\bottomrule
\end{tabular}
\end{table}

% --- Probe budget ablation ---
\paragraph{Effect of probe budget $B_{\mathrm{probe}}$.}
Table~\ref{tab:ablation-probe} fixes $K = 2$ and varies
$B_{\mathrm{probe}} \in \{128, 170, 256, 512\}$.
Larger probe budgets allow the exploration paths to develop more
complete reasoning chains, increasing both the validity rate $\nu$ and
the reliability of the consensus signal.
However, beyond $B_{\mathrm{probe}} = 256$, the additional cost is not
justified: the marginal accuracy gain is XX.X\,pp while the per-sample
cost increases by XX.X tokens.
At $B_{\mathrm{probe}} = 512$, the exploration phase alone consumes as
many tokens as a single Fixed-512 run, eliminating the efficiency
advantage of adaptive routing.

\begin{table}[t]
\centering
\caption{%
  \textbf{Ablation on $B_{\mathrm{probe}}$} with $K = 2$,
  $B_{\mathrm{med}} = 256$, $B_{\mathrm{hard}} = 512$, default thresholds.
}
\label{tab:ablation-probe}
\small
\begin{tabular}{@{}ccccccc@{}}
\toprule
$B_{\mathrm{probe}}$ & \textbf{Acc (\%)} & \textbf{Avg Tok} &
  \textbf{E\%} & \textbf{M\%} & \textbf{H\%} \\
\midrule
128 & XX.X & XX.X & XX.X & XX.X & XX.X \\
170 & XX.X & XX.X & XX.X & XX.X & XX.X \\
256 & XX.X & XX.X & XX.X & XX.X & XX.X \\
512 & XX.X & XX.X & XX.X & XX.X & XX.X \\
\bottomrule
\end{tabular}
\end{table}

% --- Threshold ablation ---
\paragraph{Sensitivity to routing thresholds.}
Table~\ref{tab:ablation-thresh} fixes $K = 4$, $B_{\mathrm{probe}} = 128$
and sweeps the easy-route threshold
$\theta_E \in \{0.75, 0.90, 0.95\}$ and medium-route threshold
$\theta_M \in \{0.50, 0.60, 0.70\}$.
Higher $\theta_E$ makes the router more conservative, sending fewer
problems to the easy route and increasing average cost.
The accuracy is relatively insensitive to $\theta_E$ in the range
$[0.75, 0.95]$, varying by only XX.X\,pp, because problems
mis-classified as easy at $\theta_E = 0.75$ are already near the
decision boundary and contribute minimal accuracy loss.
The medium threshold $\theta_M$ has a larger effect: raising $\theta_M$
from 0.50 to 0.70 shifts samples from the medium to the hard route,
increasing both accuracy (by XX.X\,pp) and cost (by XX.X tokens).

\begin{table}[t]
\centering
\caption{%
  \textbf{Ablation on routing thresholds} with $K = 4$,
  $B_{\mathrm{probe}} = 128$. Each cell: Acc (\%) / Avg Tok.
}
\label{tab:ablation-thresh}
\small
\begin{tabular}{@{}lccc@{}}
\toprule
& \multicolumn{3}{c}{$\theta_M$ (medium threshold)} \\
\cmidrule(lr){2-4}
$\theta_E$ (easy) & 0.50 & 0.60 & 0.70 \\
\midrule
0.75 & XX.X\,/\,XX.X & XX.X\,/\,XX.X & XX.X\,/\,XX.X \\
0.90 & XX.X\,/\,XX.X & XX.X\,/\,XX.X & XX.X\,/\,XX.X \\
0.95 & XX.X\,/\,XX.X & XX.X\,/\,XX.X & XX.X\,/\,XX.X \\
\bottomrule
\end{tabular}
\end{table}

% ---------------------------------------------------------------------------
\subsection{Route Analysis}
\label{sec:exp:routes}
% ---------------------------------------------------------------------------

To understand \emph{how} the consensus router allocates compute, we
analyze the behavior of each difficulty route in detail.

\paragraph{Per-route accuracy.}
Table~\ref{tab:route-accuracy} breaks down accuracy by assigned route
for the best \method{} configuration ($K = 4$, $B_{\mathrm{probe}} = 128$).
Easy-route samples achieve the highest accuracy (XX.X\%), confirming
that the consensus signal reliably identifies problems the model can
solve with minimal compute.
Medium-route accuracy (XX.X\%) is intermediate, while hard-route accuracy
(XX.X\%) is lowest---as expected, since these are the problems where the
model genuinely struggles.
Critically, the easy-route accuracy exceeds what these same samples
achieve under Fixed-512 (XX.X\%), demonstrating that extended reasoning
\emph{hurts} on easy problems (the overthinking phenomenon discussed
in \S\ref{sec:intro}).

\begin{table}[t]
\centering
\caption{%
  \textbf{Per-route analysis} for \method{} ($K = 4$,
  $B_{\mathrm{probe}} = 128$, $n = 200$).
  ``Fixed-512 acc on same'' reports the accuracy of Fixed-512 on the
  subset of samples assigned to each route.
}
\label{tab:route-accuracy}
\small
\begin{tabular}{@{}lcccc@{}}
\toprule
\textbf{Route} & \textbf{\% Samples} & \textbf{Acc (\%)} &
  \textbf{Avg Tok} & \textbf{Fixed-512 acc on same (\%)} \\
\midrule
Easy   & XX.X & XX.X & XX.X & XX.X \\
Medium & XX.X & XX.X & XX.X & XX.X \\
Hard   & XX.X & XX.X & XX.X & XX.X \\
\midrule
Overall & 100.0 & XX.X & XX.X & XX.X \\
\bottomrule
\end{tabular}
\end{table}

\paragraph{Token savings per route.}
The easy route consumes only $K \times B_{\mathrm{probe}}$ tokens
(XX.X on average), saving XX.X\% relative to Fixed-512.
The medium route consumes an average of XX.X tokens (exploration +
extension), while the hard route consumes XX.X tokens (exploration +
full deliberation).
Across all routes, the weighted average token consumption is XX.X,
representing a XX.X\% reduction compared to Fixed-512.

\paragraph{Consensus score and correctness.}
Figure~\ref{fig:consensus-correct} plots the relationship between the
consensus confidence score $\gamma$ and per-sample correctness.
Samples with $\gamma \geq 0.8$ are correct XX.X\% of the time, while
samples with $\gamma < 0.4$ are correct only XX.X\% of the time.
The Spearman rank correlation between $\gamma$ and binary correctness
is $\rho = $ XX.XX ($p < $ XX), empirically validating
Proposition~\ref{prop:monotone}'s prediction that consensus is a
monotone difficulty estimator.
This strong correlation explains why the simple three-tier routing rule
in Eq.~\eqref{eq:routing} is effective: the consensus signal separates
difficulty levels well enough that even coarse binning into three routes
captures most of the available adaptive gain.

\begin{figure}[t]
\centering
% \includegraphics[width=0.55\textwidth]{figures/consensus_vs_correctness.pdf}
\fbox{\parbox{0.55\textwidth}{\centering\vspace{2cm}%
  \textit{Placeholder: scatter/binned plot of consensus score $\gamma$
  vs.\ correctness rate, with fitted logistic curve.}%
\vspace{2cm}}}
\caption{%
  \textbf{Consensus score vs.\ correctness.}
  Each point represents a bin of samples grouped by consensus confidence
  $\gamma$; the $y$-axis shows the fraction of correct answers.
  The strong positive correlation ($\rho =$ XX.XX) confirms that
  multi-path consensus is a reliable difficulty signal.
}
\label{fig:consensus-correct}
\end{figure}


% ---------------------------------------------------------------------------
\subsection{Experimental Setup}
\label{sec:exp:setup}
% ---------------------------------------------------------------------------

\paragraph{Model and inference.}
All experiments use \textbf{Qwen3-8B} with thinking mode enabled
(\texttt{enable\_thinking=True}), which instructs the model to produce an
extended chain-of-thought within \texttt{<think>...</think>} tags before
emitting a final answer.
We control the thinking token budget $b$ by setting
\texttt{max\_thinking\_tokens}; once this limit is reached, the model is
forced to exit the thinking block and produce an answer.
Greedy decoding ($\tau = 0$) is used for the first exploration path and
all resolution phases; the remaining $K{-}1$ exploration paths use
$\tau = 0.7$ to encourage diversity (\S\ref{sec:method:explore}).

\paragraph{Benchmark.}
We evaluate on \textbf{GSM8K}~\citep{cobbe2021training}, a dataset of
1{,}319 grade-school math word problems.
For the main experiments and ablations we sample subsets of $n = 200$
problems (with fixed random seeds for reproducibility); full-dataset
results are reported for the best configuration.
Accuracy is measured via exact match after deterministic number extraction
(the rightmost numerical value in the model's output).

\paragraph{Baselines.}
We compare \method{} against four baselines that span the natural
design space:
\begin{itemize}[nosep,leftmargin=*]
  \item \textbf{Fixed-$b$}: single-pass generation with a fixed thinking
    budget $b \in \{128, 256, 512\}$ tokens.
  \item \textbf{SC@$K{\times}128$}: self-consistency~\citep{wang2023selfconsistency}
    with $K$ independent samples at $b = 128$ tokens each, followed by
    majority voting.
    This baseline uses the same total exploration tokens as \method{}
    with $K$ paths and $B_{\mathrm{probe}} = 128$, providing a
    compute-matched comparison.
\end{itemize}

\paragraph{\method{} configurations.}
We explore the following hyperparameter axes:
parallelism $K \in \{2, 3, 4, 8\}$,
probe budget $B_{\mathrm{probe}} \in \{128, 170, 256, 512\}$,
medium resolution budget $B_{\mathrm{med}} \in \{256, 512\}$, and
hard resolution budget $B_{\mathrm{hard}} \in \{512, 1024\}$.
Routing thresholds default to $\theta_E = 0.75$ (easy) and
$\theta_M = 0.5$ (medium) unless varied in ablation studies.
The consensus confidence weights are fixed at
$\gamma = 0.5\alpha + 0.3\nu + 0.2\mu$ throughout.

\paragraph{Evaluation metrics.}
For each method we report:
(1)~\textbf{Accuracy} (exact-match \%),
(2)~\textbf{Avg Tokens} (mean thinking tokens consumed per sample,
including exploration overhead for \method{}),
(3)~\textbf{Route distribution} (fraction of samples classified as
Easy / Medium / Hard by the consensus router).

\paragraph{Hardware.}
All experiments are conducted on a single NVIDIA A100-80GB GPU.

% ---------------------------------------------------------------------------
\subsection{Main Results}
\label{sec:exp:main}
% ---------------------------------------------------------------------------

Table~\ref{tab:main-rs} presents the primary comparison between
\method{}, fixed-budget baselines, and self-consistency on GSM8K
($n = 200$).
All \method{} variants use default routing thresholds
($\theta_E = 0.75$, $\theta_M = 0.5$).

\begin{table}[t]
\centering
\caption{%
  \textbf{Main results on GSM8K} ($n = 200$, Qwen3-8B thinking mode).
  \textbf{Avg Tok} includes full end-to-end cost (exploration +
  resolution).
  Route columns report the fraction of samples classified as Easy (E),
  Medium (M), and Hard (H) by the consensus router.
  Best adaptive result in \textbf{bold}; best baseline \underline{underlined}.
}
\label{tab:main-rs}
\small
\setlength{\tabcolsep}{4pt}
\begin{tabular}{@{}lcccccccc@{}}
\toprule
\textbf{Method} & $K$ & $B_{\mathrm{probe}}$ &
  \textbf{Acc (\%)} & \textbf{Avg Tok} &
  \textbf{E\%} & \textbf{M\%} & \textbf{H\%} \\
\midrule
\multicolumn{8}{@{}l}{\textit{Fixed-budget baselines}} \\
Fixed-128       & -- & -- & XX.X & 128  & -- & -- & -- \\
Fixed-256       & -- & -- & XX.X & 256  & -- & -- & -- \\
\underline{Fixed-512} & -- & -- & \underline{XX.X} & 512  & -- & -- & -- \\
\midrule
\multicolumn{8}{@{}l}{\textit{Self-consistency baselines}} \\
SC@$2{\times}128$  & 2 & 128 & XX.X & 256  & -- & -- & -- \\
SC@$4{\times}128$  & 4 & 128 & XX.X & 512  & -- & -- & -- \\
SC@$8{\times}128$  & 8 & 128 & XX.X & 1024 & -- & -- & -- \\
\midrule
\multicolumn{8}{@{}l}{\textit{\method{} (ours)}} \\
\method{}-2      & 2 & 128 & XX.X & XX.X & XX.X & XX.X & XX.X \\
\method{}-4      & 4 & 128 & XX.X & XX.X & XX.X & XX.X & XX.X \\
\method{}-4-256  & 4 & 256 & XX.X & XX.X & XX.X & XX.X & XX.X \\
\textbf{\method{}-8} & 8 & 128 & \textbf{XX.X} & XX.X & XX.X & XX.X & XX.X \\
\method{}-8-256  & 8 & 256 & XX.X & XX.X & XX.X & XX.X & XX.X \\
\bottomrule
\end{tabular}
\end{table}

Several observations emerge from Table~\ref{tab:main-rs}.
First, \method{} consistently outperforms fixed-budget baselines at
comparable or lower token cost.
The best configuration (\method{}-8) achieves XX.X\% accuracy using an
average of XX.X tokens---a \textbf{+XX.X\,pp} improvement over Fixed-512
at \textbf{XX\%} of the token budget.
Second, \method{} substantially outperforms the compute-matched
self-consistency baseline (SC@$K{\times}128$), confirming that the
consensus signal is more effectively leveraged through adaptive routing
than through simple majority voting.
For instance, at $K = 4$, \method{}-4 achieves XX.X\% vs.\
SC@$4{\times}128$'s XX.X\%, while consuming fewer total tokens because
easy samples are resolved without a resolution phase.
Third, the route distribution reveals that a substantial fraction of
GSM8K problems (XX.X\%) are routed as easy, validating the central
premise that many problems can be solved with minimal compute.

% ---------------------------------------------------------------------------
\subsection{Ablation Studies}
\label{sec:exp:ablation}
% ---------------------------------------------------------------------------

We systematically ablate the three key hyperparameters of \method{} to
understand their individual contributions.

% --- K ablation ---
\paragraph{Effect of parallelism $K$.}
Table~\ref{tab:ablation-k} fixes $B_{\mathrm{probe}} = 128$ and varies
$K \in \{2, 3, 4, 8\}$.
Increasing $K$ improves the reliability of the consensus signal
(Proposition~\ref{prop:monotone}), yielding monotonically higher accuracy.
However, the exploration cost scales linearly in $K$, creating a
diminishing-returns trade-off: from $K = 4$ to $K = 8$, accuracy improves
by XX.X\,pp while average tokens increase by XX.X.
We find $K = 4$ to be the best cost-accuracy operating point for GSM8K.

\begin{table}[t]
\centering
\caption{%
  \textbf{Ablation on $K$} (parallelism) with $B_{\mathrm{probe}} = 128$,
  $B_{\mathrm{med}} = 256$, $B_{\mathrm{hard}} = 512$, default thresholds.
}
\label{tab:ablation-k}
\small
\begin{tabular}{@{}ccccccc@{}}
\toprule
$K$ & \textbf{Acc (\%)} & \textbf{Avg Tok} &
  \textbf{E\%} & \textbf{M\%} & \textbf{H\%} \\
\midrule
2 & XX.X & XX.X & XX.X & XX.X & XX.X \\
3 & XX.X & XX.X & XX.X & XX.X & XX.X \\
4 & XX.X & XX.X & XX.X & XX.X & XX.X \\
8 & XX.X & XX.X & XX.X & XX.X & XX.X \\
\bottomrule
\end{tabular}
\end{table}

% --- Probe budget ablation ---
\paragraph{Effect of probe budget $B_{\mathrm{probe}}$.}
Table~\ref{tab:ablation-probe} fixes $K = 2$ and varies
$B_{\mathrm{probe}} \in \{128, 170, 256, 512\}$.
Larger probe budgets allow the exploration paths to develop more
complete reasoning chains, increasing both the validity rate $\nu$ and
the reliability of the consensus signal.
However, beyond $B_{\mathrm{probe}} = 256$, the additional cost is not
justified: the marginal accuracy gain is XX.X\,pp while the per-sample
cost increases by XX.X tokens.
At $B_{\mathrm{probe}} = 512$, the exploration phase alone consumes as
many tokens as a single Fixed-512 run, eliminating the efficiency
advantage of adaptive routing.

\begin{table}[t]
\centering
\caption{%
  \textbf{Ablation on $B_{\mathrm{probe}}$} with $K = 2$,
  $B_{\mathrm{med}} = 256$, $B_{\mathrm{hard}} = 512$, default thresholds.
}
\label{tab:ablation-probe}
\small
\begin{tabular}{@{}ccccccc@{}}
\toprule
$B_{\mathrm{probe}}$ & \textbf{Acc (\%)} & \textbf{Avg Tok} &
  \textbf{E\%} & \textbf{M\%} & \textbf{H\%} \\
\midrule
128 & XX.X & XX.X & XX.X & XX.X & XX.X \\
170 & XX.X & XX.X & XX.X & XX.X & XX.X \\
256 & XX.X & XX.X & XX.X & XX.X & XX.X \\
512 & XX.X & XX.X & XX.X & XX.X & XX.X \\
\bottomrule
\end{tabular}
\end{table}

% --- Threshold ablation ---
\paragraph{Sensitivity to routing thresholds.}
Table~\ref{tab:ablation-thresh} fixes $K = 4$, $B_{\mathrm{probe}} = 128$
and sweeps the easy-route threshold
$\theta_E \in \{0.75, 0.90, 0.95\}$ and medium-route threshold
$\theta_M \in \{0.50, 0.60, 0.70\}$.
Higher $\theta_E$ makes the router more conservative, sending fewer
problems to the easy route and increasing average cost.
The accuracy is relatively insensitive to $\theta_E$ in the range
$[0.75, 0.95]$, varying by only XX.X\,pp, because problems
mis-classified as easy at $\theta_E = 0.75$ are already near the
decision boundary and contribute minimal accuracy loss.
The medium threshold $\theta_M$ has a larger effect: raising $\theta_M$
from 0.50 to 0.70 shifts samples from the medium to the hard route,
increasing both accuracy (by XX.X\,pp) and cost (by XX.X tokens).

\begin{table}[t]
\centering
\caption{%
  \textbf{Ablation on routing thresholds} with $K = 4$,
  $B_{\mathrm{probe}} = 128$. Each cell: Acc (\%) / Avg Tok.
}
\label{tab:ablation-thresh}
\small
\begin{tabular}{@{}lccc@{}}
\toprule
& \multicolumn{3}{c}{$\theta_M$ (medium threshold)} \\
\cmidrule(lr){2-4}
$\theta_E$ (easy) & 0.50 & 0.60 & 0.70 \\
\midrule
0.75 & XX.X\,/\,XX.X & XX.X\,/\,XX.X & XX.X\,/\,XX.X \\
0.90 & XX.X\,/\,XX.X & XX.X\,/\,XX.X & XX.X\,/\,XX.X \\
0.95 & XX.X\,/\,XX.X & XX.X\,/\,XX.X & XX.X\,/\,XX.X \\
\bottomrule
\end{tabular}
\end{table}

% ---------------------------------------------------------------------------
\subsection{Route Analysis}
\label{sec:exp:routes}
% ---------------------------------------------------------------------------

To understand \emph{how} the consensus router allocates compute, we
analyze the behavior of each difficulty route in detail.

\paragraph{Per-route accuracy.}
Table~\ref{tab:route-accuracy} breaks down accuracy by assigned route
for the best \method{} configuration ($K = 4$, $B_{\mathrm{probe}} = 128$).
Easy-route samples achieve the highest accuracy (XX.X\%), confirming
that the consensus signal reliably identifies problems the model can
solve with minimal compute.
Medium-route accuracy (XX.X\%) is intermediate, while hard-route accuracy
(XX.X\%) is lowest---as expected, since these are the problems where the
model genuinely struggles.
Critically, the easy-route accuracy exceeds what these same samples
achieve under Fixed-512 (XX.X\%), demonstrating that extended reasoning
\emph{hurts} on easy problems (the overthinking phenomenon discussed
in \S\ref{sec:intro}).

\begin{table}[t]
\centering
\caption{%
  \textbf{Per-route analysis} for \method{} ($K = 4$,
  $B_{\mathrm{probe}} = 128$, $n = 200$).
  ``Fixed-512 acc on same'' reports the accuracy of Fixed-512 on the
  subset of samples assigned to each route.
}
\label{tab:route-accuracy}
\small
\begin{tabular}{@{}lcccc@{}}
\toprule
\textbf{Route} & \textbf{\% Samples} & \textbf{Acc (\%)} &
  \textbf{Avg Tok} & \textbf{Fixed-512 acc on same (\%)} \\
\midrule
Easy   & XX.X & XX.X & XX.X & XX.X \\
Medium & XX.X & XX.X & XX.X & XX.X \\
Hard   & XX.X & XX.X & XX.X & XX.X \\
\midrule
Overall & 100.0 & XX.X & XX.X & XX.X \\
\bottomrule
\end{tabular}
\end{table}

\paragraph{Token savings per route.}
The easy route consumes only $K \times B_{\mathrm{probe}}$ tokens
(XX.X on average), saving XX.X\% relative to Fixed-512.
The medium route consumes an average of XX.X tokens (exploration +
extension), while the hard route consumes XX.X tokens (exploration +
full deliberation).
Across all routes, the weighted average token consumption is XX.X,
representing a XX.X\% reduction compared to Fixed-512.

\paragraph{Consensus score and correctness.}
Figure~\ref{fig:consensus-correct} plots the relationship between the
consensus confidence score $\gamma$ and per-sample correctness.
Samples with $\gamma \geq 0.8$ are correct XX.X\% of the time, while
samples with $\gamma < 0.4$ are correct only XX.X\% of the time.
The Spearman rank correlation between $\gamma$ and binary correctness
is $\rho = $ XX.XX ($p < $ XX), empirically validating
Proposition~\ref{prop:monotone}'s prediction that consensus is a
monotone difficulty estimator.
This strong correlation explains why the simple three-tier routing rule
in Eq.~\eqref{eq:routing} is effective: the consensus signal separates
difficulty levels well enough that even coarse binning into three routes
captures most of the available adaptive gain.

\begin{figure}[t]
\centering
% \includegraphics[width=0.55\textwidth]{figures/consensus_vs_correctness.pdf}
\fbox{\parbox{0.55\textwidth}{\centering\vspace{2cm}%
  \textit{Placeholder: scatter/binned plot of consensus score $\gamma$
  vs.\ correctness rate, with fitted logistic curve.}%
\vspace{2cm}}}
\caption{%
  \textbf{Consensus score vs.\ correctness.}
  Each point represents a bin of samples grouped by consensus confidence
  $\gamma$; the $y$-axis shows the fraction of correct answers.
  The strong positive correlation ($\rho =$ XX.XX) confirms that
  multi-path consensus is a reliable difficulty signal.
}
\label{fig:consensus-correct}
\end{figure}


% ---------------------------------------------------------------------------
\subsection{Experimental Setup}
\label{sec:exp:setup}
% ---------------------------------------------------------------------------

\paragraph{Model and inference.}
All experiments use \textbf{Qwen3-8B} with thinking mode enabled
(\texttt{enable\_thinking=True}), which instructs the model to produce an
extended chain-of-thought within \texttt{<think>...</think>} tags before
emitting a final answer.
We control the thinking token budget $b$ by setting
\texttt{max\_thinking\_tokens}; once this limit is reached, the model is
forced to exit the thinking block and produce an answer.
Greedy decoding ($\tau = 0$) is used for the first exploration path and
all resolution phases; the remaining $K{-}1$ exploration paths use
$\tau = 0.7$ to encourage diversity (\S\ref{sec:method:explore}).

\paragraph{Benchmark.}
We evaluate on \textbf{GSM8K}~\citep{cobbe2021training}, a dataset of
1{,}319 grade-school math word problems.
For the main experiments and ablations we sample subsets of $n = 200$
problems (with fixed random seeds for reproducibility); full-dataset
results are reported for the best configuration.
Accuracy is measured via exact match after deterministic number extraction
(the rightmost numerical value in the model's output).

\paragraph{Baselines.}
We compare \method{} against four baselines that span the natural
design space:
\begin{itemize}[nosep,leftmargin=*]
  \item \textbf{Fixed-$b$}: single-pass generation with a fixed thinking
    budget $b \in \{128, 256, 512\}$ tokens.
  \item \textbf{SC@$K{\times}128$}: self-consistency~\citep{wang2023selfconsistency}
    with $K$ independent samples at $b = 128$ tokens each, followed by
    majority voting.
    This baseline uses the same total exploration tokens as \method{}
    with $K$ paths and $B_{\mathrm{probe}} = 128$, providing a
    compute-matched comparison.
\end{itemize}

\paragraph{\method{} configurations.}
We explore the following hyperparameter axes:
parallelism $K \in \{2, 3, 4, 8\}$,
probe budget $B_{\mathrm{probe}} \in \{128, 170, 256, 512\}$,
medium resolution budget $B_{\mathrm{med}} \in \{256, 512\}$, and
hard resolution budget $B_{\mathrm{hard}} \in \{512, 1024\}$.
Routing thresholds default to $\theta_E = 0.75$ (easy) and
$\theta_M = 0.5$ (medium) unless varied in ablation studies.
The consensus confidence weights are fixed at
$\gamma = 0.5\alpha + 0.3\nu + 0.2\mu$ throughout.

\paragraph{Evaluation metrics.}
For each method we report:
(1)~\textbf{Accuracy} (exact-match \%),
(2)~\textbf{Avg Tokens} (mean thinking tokens consumed per sample,
including exploration overhead for \method{}),
(3)~\textbf{Route distribution} (fraction of samples classified as
Easy / Medium / Hard by the consensus router).

\paragraph{Hardware.}
All experiments are conducted on a single NVIDIA A100-80GB GPU.

% ---------------------------------------------------------------------------
\subsection{Main Results}
\label{sec:exp:main}
% ---------------------------------------------------------------------------

Table~\ref{tab:main-rs} presents the primary comparison between
\method{}, fixed-budget baselines, and self-consistency on GSM8K
($n = 200$).
All \method{} variants use default routing thresholds
($\theta_E = 0.75$, $\theta_M = 0.5$).

\begin{table}[t]
\centering
\caption{%
  \textbf{Main results on GSM8K} ($n = 200$, Qwen3-8B thinking mode).
  \textbf{Avg Tok} includes full end-to-end cost (exploration +
  resolution).
  Route columns report the fraction of samples classified as Easy (E),
  Medium (M), and Hard (H) by the consensus router.
  Best adaptive result in \textbf{bold}; best baseline \underline{underlined}.
}
\label{tab:main-rs}
\small
\setlength{\tabcolsep}{4pt}
\begin{tabular}{@{}lcccccccc@{}}
\toprule
\textbf{Method} & $K$ & $B_{\mathrm{probe}}$ &
  \textbf{Acc (\%)} & \textbf{Avg Tok} &
  \textbf{E\%} & \textbf{M\%} & \textbf{H\%} \\
\midrule
\multicolumn{8}{@{}l}{\textit{Fixed-budget baselines}} \\
Fixed-128       & -- & -- & XX.X & 128  & -- & -- & -- \\
Fixed-256       & -- & -- & XX.X & 256  & -- & -- & -- \\
\underline{Fixed-512} & -- & -- & \underline{XX.X} & 512  & -- & -- & -- \\
\midrule
\multicolumn{8}{@{}l}{\textit{Self-consistency baselines}} \\
SC@$2{\times}128$  & 2 & 128 & XX.X & 256  & -- & -- & -- \\
SC@$4{\times}128$  & 4 & 128 & XX.X & 512  & -- & -- & -- \\
SC@$8{\times}128$  & 8 & 128 & XX.X & 1024 & -- & -- & -- \\
\midrule
\multicolumn{8}{@{}l}{\textit{\method{} (ours)}} \\
\method{}-2      & 2 & 128 & XX.X & XX.X & XX.X & XX.X & XX.X \\
\method{}-4      & 4 & 128 & XX.X & XX.X & XX.X & XX.X & XX.X \\
\method{}-4-256  & 4 & 256 & XX.X & XX.X & XX.X & XX.X & XX.X \\
\textbf{\method{}-8} & 8 & 128 & \textbf{XX.X} & XX.X & XX.X & XX.X & XX.X \\
\method{}-8-256  & 8 & 256 & XX.X & XX.X & XX.X & XX.X & XX.X \\
\bottomrule
\end{tabular}
\end{table}

Several observations emerge from Table~\ref{tab:main-rs}.
First, \method{} consistently outperforms fixed-budget baselines at
comparable or lower token cost.
The best configuration (\method{}-8) achieves XX.X\% accuracy using an
average of XX.X tokens---a \textbf{+XX.X\,pp} improvement over Fixed-512
at \textbf{XX\%} of the token budget.
Second, \method{} substantially outperforms the compute-matched
self-consistency baseline (SC@$K{\times}128$), confirming that the
consensus signal is more effectively leveraged through adaptive routing
than through simple majority voting.
For instance, at $K = 4$, \method{}-4 achieves XX.X\% vs.\
SC@$4{\times}128$'s XX.X\%, while consuming fewer total tokens because
easy samples are resolved without a resolution phase.
Third, the route distribution reveals that a substantial fraction of
GSM8K problems (XX.X\%) are routed as easy, validating the central
premise that many problems can be solved with minimal compute.

% ---------------------------------------------------------------------------
\subsection{Ablation Studies}
\label{sec:exp:ablation}
% ---------------------------------------------------------------------------

We systematically ablate the three key hyperparameters of \method{} to
understand their individual contributions.

% --- K ablation ---
\paragraph{Effect of parallelism $K$.}
Table~\ref{tab:ablation-k} fixes $B_{\mathrm{probe}} = 128$ and varies
$K \in \{2, 3, 4, 8\}$.
Increasing $K$ improves the reliability of the consensus signal
(Proposition~\ref{prop:monotone}), yielding monotonically higher accuracy.
However, the exploration cost scales linearly in $K$, creating a
diminishing-returns trade-off: from $K = 4$ to $K = 8$, accuracy improves
by XX.X\,pp while average tokens increase by XX.X.
We find $K = 4$ to be the best cost-accuracy operating point for GSM8K.

\begin{table}[t]
\centering
\caption{%
  \textbf{Ablation on $K$} (parallelism) with $B_{\mathrm{probe}} = 128$,
  $B_{\mathrm{med}} = 256$, $B_{\mathrm{hard}} = 512$, default thresholds.
}
\label{tab:ablation-k}
\small
\begin{tabular}{@{}ccccccc@{}}
\toprule
$K$ & \textbf{Acc (\%)} & \textbf{Avg Tok} &
  \textbf{E\%} & \textbf{M\%} & \textbf{H\%} \\
\midrule
2 & XX.X & XX.X & XX.X & XX.X & XX.X \\
3 & XX.X & XX.X & XX.X & XX.X & XX.X \\
4 & XX.X & XX.X & XX.X & XX.X & XX.X \\
8 & XX.X & XX.X & XX.X & XX.X & XX.X \\
\bottomrule
\end{tabular}
\end{table}

% --- Probe budget ablation ---
\paragraph{Effect of probe budget $B_{\mathrm{probe}}$.}
Table~\ref{tab:ablation-probe} fixes $K = 2$ and varies
$B_{\mathrm{probe}} \in \{128, 170, 256, 512\}$.
Larger probe budgets allow the exploration paths to develop more
complete reasoning chains, increasing both the validity rate $\nu$ and
the reliability of the consensus signal.
However, beyond $B_{\mathrm{probe}} = 256$, the additional cost is not
justified: the marginal accuracy gain is XX.X\,pp while the per-sample
cost increases by XX.X tokens.
At $B_{\mathrm{probe}} = 512$, the exploration phase alone consumes as
many tokens as a single Fixed-512 run, eliminating the efficiency
advantage of adaptive routing.

\begin{table}[t]
\centering
\caption{%
  \textbf{Ablation on $B_{\mathrm{probe}}$} with $K = 2$,
  $B_{\mathrm{med}} = 256$, $B_{\mathrm{hard}} = 512$, default thresholds.
}
\label{tab:ablation-probe}
\small
\begin{tabular}{@{}ccccccc@{}}
\toprule
$B_{\mathrm{probe}}$ & \textbf{Acc (\%)} & \textbf{Avg Tok} &
  \textbf{E\%} & \textbf{M\%} & \textbf{H\%} \\
\midrule
128 & XX.X & XX.X & XX.X & XX.X & XX.X \\
170 & XX.X & XX.X & XX.X & XX.X & XX.X \\
256 & XX.X & XX.X & XX.X & XX.X & XX.X \\
512 & XX.X & XX.X & XX.X & XX.X & XX.X \\
\bottomrule
\end{tabular}
\end{table}

% --- Threshold ablation ---
\paragraph{Sensitivity to routing thresholds.}
Table~\ref{tab:ablation-thresh} fixes $K = 4$, $B_{\mathrm{probe}} = 128$
and sweeps the easy-route threshold
$\theta_E \in \{0.75, 0.90, 0.95\}$ and medium-route threshold
$\theta_M \in \{0.50, 0.60, 0.70\}$.
Higher $\theta_E$ makes the router more conservative, sending fewer
problems to the easy route and increasing average cost.
The accuracy is relatively insensitive to $\theta_E$ in the range
$[0.75, 0.95]$, varying by only XX.X\,pp, because problems
mis-classified as easy at $\theta_E = 0.75$ are already near the
decision boundary and contribute minimal accuracy loss.
The medium threshold $\theta_M$ has a larger effect: raising $\theta_M$
from 0.50 to 0.70 shifts samples from the medium to the hard route,
increasing both accuracy (by XX.X\,pp) and cost (by XX.X tokens).

\begin{table}[t]
\centering
\caption{%
  \textbf{Ablation on routing thresholds} with $K = 4$,
  $B_{\mathrm{probe}} = 128$. Each cell: Acc (\%) / Avg Tok.
}
\label{tab:ablation-thresh}
\small
\begin{tabular}{@{}lccc@{}}
\toprule
& \multicolumn{3}{c}{$\theta_M$ (medium threshold)} \\
\cmidrule(lr){2-4}
$\theta_E$ (easy) & 0.50 & 0.60 & 0.70 \\
\midrule
0.75 & XX.X\,/\,XX.X & XX.X\,/\,XX.X & XX.X\,/\,XX.X \\
0.90 & XX.X\,/\,XX.X & XX.X\,/\,XX.X & XX.X\,/\,XX.X \\
0.95 & XX.X\,/\,XX.X & XX.X\,/\,XX.X & XX.X\,/\,XX.X \\
\bottomrule
\end{tabular}
\end{table}

% ---------------------------------------------------------------------------
\subsection{Route Analysis}
\label{sec:exp:routes}
% ---------------------------------------------------------------------------

To understand \emph{how} the consensus router allocates compute, we
analyze the behavior of each difficulty route in detail.

\paragraph{Per-route accuracy.}
Table~\ref{tab:route-accuracy} breaks down accuracy by assigned route
for the best \method{} configuration ($K = 4$, $B_{\mathrm{probe}} = 128$).
Easy-route samples achieve the highest accuracy (XX.X\%), confirming
that the consensus signal reliably identifies problems the model can
solve with minimal compute.
Medium-route accuracy (XX.X\%) is intermediate, while hard-route accuracy
(XX.X\%) is lowest---as expected, since these are the problems where the
model genuinely struggles.
Critically, the easy-route accuracy exceeds what these same samples
achieve under Fixed-512 (XX.X\%), demonstrating that extended reasoning
\emph{hurts} on easy problems (the overthinking phenomenon discussed
in \S\ref{sec:intro}).

\begin{table}[t]
\centering
\caption{%
  \textbf{Per-route analysis} for \method{} ($K = 4$,
  $B_{\mathrm{probe}} = 128$, $n = 200$).
  ``Fixed-512 acc on same'' reports the accuracy of Fixed-512 on the
  subset of samples assigned to each route.
}
\label{tab:route-accuracy}
\small
\begin{tabular}{@{}lcccc@{}}
\toprule
\textbf{Route} & \textbf{\% Samples} & \textbf{Acc (\%)} &
  \textbf{Avg Tok} & \textbf{Fixed-512 acc on same (\%)} \\
\midrule
Easy   & XX.X & XX.X & XX.X & XX.X \\
Medium & XX.X & XX.X & XX.X & XX.X \\
Hard   & XX.X & XX.X & XX.X & XX.X \\
\midrule
Overall & 100.0 & XX.X & XX.X & XX.X \\
\bottomrule
\end{tabular}
\end{table}

\paragraph{Token savings per route.}
The easy route consumes only $K \times B_{\mathrm{probe}}$ tokens
(XX.X on average), saving XX.X\% relative to Fixed-512.
The medium route consumes an average of XX.X tokens (exploration +
extension), while the hard route consumes XX.X tokens (exploration +
full deliberation).
Across all routes, the weighted average token consumption is XX.X,
representing a XX.X\% reduction compared to Fixed-512.

\paragraph{Consensus score and correctness.}
Figure~\ref{fig:consensus-correct} plots the relationship between the
consensus confidence score $\gamma$ and per-sample correctness.
Samples with $\gamma \geq 0.8$ are correct XX.X\% of the time, while
samples with $\gamma < 0.4$ are correct only XX.X\% of the time.
The Spearman rank correlation between $\gamma$ and binary correctness
is $\rho = $ XX.XX ($p < $ XX), empirically validating
Proposition~\ref{prop:monotone}'s prediction that consensus is a
monotone difficulty estimator.
This strong correlation explains why the simple three-tier routing rule
in Eq.~\eqref{eq:routing} is effective: the consensus signal separates
difficulty levels well enough that even coarse binning into three routes
captures most of the available adaptive gain.

\begin{figure}[t]
\centering
% \includegraphics[width=0.55\textwidth]{figures/consensus_vs_correctness.pdf}
\fbox{\parbox{0.55\textwidth}{\centering\vspace{2cm}%
  \textit{Placeholder: scatter/binned plot of consensus score $\gamma$
  vs.\ correctness rate, with fitted logistic curve.}%
\vspace{2cm}}}
\caption{%
  \textbf{Consensus score vs.\ correctness.}
  Each point represents a bin of samples grouped by consensus confidence
  $\gamma$; the $y$-axis shows the fraction of correct answers.
  The strong positive correlation ($\rho =$ XX.XX) confirms that
  multi-path consensus is a reliable difficulty signal.
}
\label{fig:consensus-correct}
\end{figure}


% ---------------------------------------------------------------------------
\subsection{Experimental Setup}
\label{sec:exp:setup}
% ---------------------------------------------------------------------------

\paragraph{Model and inference.}
All experiments use \textbf{Qwen3-8B} with thinking mode enabled
(\texttt{enable\_thinking=True}), which instructs the model to produce an
extended chain-of-thought within \texttt{<think>...</think>} tags before
emitting a final answer.
We control the thinking token budget $b$ by setting
\texttt{max\_thinking\_tokens}; once this limit is reached, the model is
forced to exit the thinking block and produce an answer.
Greedy decoding ($\tau = 0$) is used for the first exploration path and
all resolution phases; the remaining $K{-}1$ exploration paths use
$\tau = 0.7$ to encourage diversity (\S\ref{sec:method:explore}).

\paragraph{Benchmark.}
We evaluate on \textbf{GSM8K}~\citep{cobbe2021training}, a dataset of
1{,}319 grade-school math word problems.
For the main experiments and ablations we sample subsets of $n = 200$
problems (with fixed random seeds for reproducibility); full-dataset
results are reported for the best configuration.
Accuracy is measured via exact match after deterministic number extraction
(the rightmost numerical value in the model's output).

\paragraph{Baselines.}
We compare \method{} against four baselines that span the natural
design space:
\begin{itemize}[nosep,leftmargin=*]
  \item \textbf{Fixed-$b$}: single-pass generation with a fixed thinking
    budget $b \in \{128, 256, 512\}$ tokens.
  \item \textbf{SC@$K{\times}128$}: self-consistency~\citep{wang2023selfconsistency}
    with $K$ independent samples at $b = 128$ tokens each, followed by
    majority voting.
    This baseline uses the same total exploration tokens as \method{}
    with $K$ paths and $B_{\mathrm{probe}} = 128$, providing a
    compute-matched comparison.
\end{itemize}

\paragraph{\method{} configurations.}
We explore the following hyperparameter axes:
parallelism $K \in \{2, 3, 4, 8\}$,
probe budget $B_{\mathrm{probe}} \in \{128, 170, 256, 512\}$,
medium resolution budget $B_{\mathrm{med}} \in \{256, 512\}$, and
hard resolution budget $B_{\mathrm{hard}} \in \{512, 1024\}$.
Routing thresholds default to $\theta_E = 0.75$ (easy) and
$\theta_M = 0.5$ (medium) unless varied in ablation studies.
The consensus confidence weights are fixed at
$\gamma = 0.5\alpha + 0.3\nu + 0.2\mu$ throughout.

\paragraph{Evaluation metrics.}
For each method we report:
(1)~\textbf{Accuracy} (exact-match \%),
(2)~\textbf{Avg Tokens} (mean thinking tokens consumed per sample,
including exploration overhead for \method{}),
(3)~\textbf{Route distribution} (fraction of samples classified as
Easy / Medium / Hard by the consensus router).

\paragraph{Hardware.}
All experiments are conducted on a single NVIDIA A100-80GB GPU.

% ---------------------------------------------------------------------------
\subsection{Main Results}
\label{sec:exp:main}
% ---------------------------------------------------------------------------

Table~\ref{tab:main-rs} presents the primary comparison between
\method{}, fixed-budget baselines, and self-consistency on GSM8K
($n = 200$).
All \method{} variants use default routing thresholds
($\theta_E = 0.75$, $\theta_M = 0.5$).

\begin{table}[t]
\centering
\caption{%
  \textbf{Main results on GSM8K} ($n = 200$, Qwen3-8B thinking mode).
  \textbf{Avg Tok} includes full end-to-end cost (exploration +
  resolution).
  Route columns report the fraction of samples classified as Easy (E),
  Medium (M), and Hard (H) by the consensus router.
  Best adaptive result in \textbf{bold}; best baseline \underline{underlined}.
}
\label{tab:main-rs}
\small
\setlength{\tabcolsep}{4pt}
\begin{tabular}{@{}lcccccccc@{}}
\toprule
\textbf{Method} & $K$ & $B_{\mathrm{probe}}$ &
  \textbf{Acc (\%)} & \textbf{Avg Tok} &
  \textbf{E\%} & \textbf{M\%} & \textbf{H\%} \\
\midrule
\multicolumn{8}{@{}l}{\textit{Fixed-budget baselines}} \\
Fixed-128       & -- & -- & XX.X & 128  & -- & -- & -- \\
Fixed-256       & -- & -- & XX.X & 256  & -- & -- & -- \\
\underline{Fixed-512} & -- & -- & \underline{XX.X} & 512  & -- & -- & -- \\
\midrule
\multicolumn{8}{@{}l}{\textit{Self-consistency baselines}} \\
SC@$2{\times}128$  & 2 & 128 & XX.X & 256  & -- & -- & -- \\
SC@$4{\times}128$  & 4 & 128 & XX.X & 512  & -- & -- & -- \\
SC@$8{\times}128$  & 8 & 128 & XX.X & 1024 & -- & -- & -- \\
\midrule
\multicolumn{8}{@{}l}{\textit{\method{} (ours)}} \\
\method{}-2      & 2 & 128 & XX.X & XX.X & XX.X & XX.X & XX.X \\
\method{}-4      & 4 & 128 & XX.X & XX.X & XX.X & XX.X & XX.X \\
\method{}-4-256  & 4 & 256 & XX.X & XX.X & XX.X & XX.X & XX.X \\
\textbf{\method{}-8} & 8 & 128 & \textbf{XX.X} & XX.X & XX.X & XX.X & XX.X \\
\method{}-8-256  & 8 & 256 & XX.X & XX.X & XX.X & XX.X & XX.X \\
\bottomrule
\end{tabular}
\end{table}

Several observations emerge from Table~\ref{tab:main-rs}.
First, \method{} consistently outperforms fixed-budget baselines at
comparable or lower token cost.
The best configuration (\method{}-8) achieves XX.X\% accuracy using an
average of XX.X tokens---a \textbf{+XX.X\,pp} improvement over Fixed-512
at \textbf{XX\%} of the token budget.
Second, \method{} substantially outperforms the compute-matched
self-consistency baseline (SC@$K{\times}128$), confirming that the
consensus signal is more effectively leveraged through adaptive routing
than through simple majority voting.
For instance, at $K = 4$, \method{}-4 achieves XX.X\% vs.\
SC@$4{\times}128$'s XX.X\%, while consuming fewer total tokens because
easy samples are resolved without a resolution phase.
Third, the route distribution reveals that a substantial fraction of
GSM8K problems (XX.X\%) are routed as easy, validating the central
premise that many problems can be solved with minimal compute.

% ---------------------------------------------------------------------------
\subsection{Ablation Studies}
\label{sec:exp:ablation}
% ---------------------------------------------------------------------------

We systematically ablate the three key hyperparameters of \method{} to
understand their individual contributions.

% --- K ablation ---
\paragraph{Effect of parallelism $K$.}
Table~\ref{tab:ablation-k} fixes $B_{\mathrm{probe}} = 128$ and varies
$K \in \{2, 3, 4, 8\}$.
Increasing $K$ improves the reliability of the consensus signal
(Proposition~\ref{prop:monotone}), yielding monotonically higher accuracy.
However, the exploration cost scales linearly in $K$, creating a
diminishing-returns trade-off: from $K = 4$ to $K = 8$, accuracy improves
by XX.X\,pp while average tokens increase by XX.X.
We find $K = 4$ to be the best cost-accuracy operating point for GSM8K.

\begin{table}[t]
\centering
\caption{%
  \textbf{Ablation on $K$} (parallelism) with $B_{\mathrm{probe}} = 128$,
  $B_{\mathrm{med}} = 256$, $B_{\mathrm{hard}} = 512$, default thresholds.
}
\label{tab:ablation-k}
\small
\begin{tabular}{@{}ccccccc@{}}
\toprule
$K$ & \textbf{Acc (\%)} & \textbf{Avg Tok} &
  \textbf{E\%} & \textbf{M\%} & \textbf{H\%} \\
\midrule
2 & XX.X & XX.X & XX.X & XX.X & XX.X \\
3 & XX.X & XX.X & XX.X & XX.X & XX.X \\
4 & XX.X & XX.X & XX.X & XX.X & XX.X \\
8 & XX.X & XX.X & XX.X & XX.X & XX.X \\
\bottomrule
\end{tabular}
\end{table}

% --- Probe budget ablation ---
\paragraph{Effect of probe budget $B_{\mathrm{probe}}$.}
Table~\ref{tab:ablation-probe} fixes $K = 2$ and varies
$B_{\mathrm{probe}} \in \{128, 170, 256, 512\}$.
Larger probe budgets allow the exploration paths to develop more
complete reasoning chains, increasing both the validity rate $\nu$ and
the reliability of the consensus signal.
However, beyond $B_{\mathrm{probe}} = 256$, the additional cost is not
justified: the marginal accuracy gain is XX.X\,pp while the per-sample
cost increases by XX.X tokens.
At $B_{\mathrm{probe}} = 512$, the exploration phase alone consumes as
many tokens as a single Fixed-512 run, eliminating the efficiency
advantage of adaptive routing.

\begin{table}[t]
\centering
\caption{%
  \textbf{Ablation on $B_{\mathrm{probe}}$} with $K = 2$,
  $B_{\mathrm{med}} = 256$, $B_{\mathrm{hard}} = 512$, default thresholds.
}
\label{tab:ablation-probe}
\small
\begin{tabular}{@{}ccccccc@{}}
\toprule
$B_{\mathrm{probe}}$ & \textbf{Acc (\%)} & \textbf{Avg Tok} &
  \textbf{E\%} & \textbf{M\%} & \textbf{H\%} \\
\midrule
128 & XX.X & XX.X & XX.X & XX.X & XX.X \\
170 & XX.X & XX.X & XX.X & XX.X & XX.X \\
256 & XX.X & XX.X & XX.X & XX.X & XX.X \\
512 & XX.X & XX.X & XX.X & XX.X & XX.X \\
\bottomrule
\end{tabular}
\end{table}

% --- Threshold ablation ---
\paragraph{Sensitivity to routing thresholds.}
Table~\ref{tab:ablation-thresh} fixes $K = 4$, $B_{\mathrm{probe}} = 128$
and sweeps the easy-route threshold
$\theta_E \in \{0.75, 0.90, 0.95\}$ and medium-route threshold
$\theta_M \in \{0.50, 0.60, 0.70\}$.
Higher $\theta_E$ makes the router more conservative, sending fewer
problems to the easy route and increasing average cost.
The accuracy is relatively insensitive to $\theta_E$ in the range
$[0.75, 0.95]$, varying by only XX.X\,pp, because problems
mis-classified as easy at $\theta_E = 0.75$ are already near the
decision boundary and contribute minimal accuracy loss.
The medium threshold $\theta_M$ has a larger effect: raising $\theta_M$
from 0.50 to 0.70 shifts samples from the medium to the hard route,
increasing both accuracy (by XX.X\,pp) and cost (by XX.X tokens).

\begin{table}[t]
\centering
\caption{%
  \textbf{Ablation on routing thresholds} with $K = 4$,
  $B_{\mathrm{probe}} = 128$. Each cell: Acc (\%) / Avg Tok.
}
\label{tab:ablation-thresh}
\small
\begin{tabular}{@{}lccc@{}}
\toprule
& \multicolumn{3}{c}{$\theta_M$ (medium threshold)} \\
\cmidrule(lr){2-4}
$\theta_E$ (easy) & 0.50 & 0.60 & 0.70 \\
\midrule
0.75 & XX.X\,/\,XX.X & XX.X\,/\,XX.X & XX.X\,/\,XX.X \\
0.90 & XX.X\,/\,XX.X & XX.X\,/\,XX.X & XX.X\,/\,XX.X \\
0.95 & XX.X\,/\,XX.X & XX.X\,/\,XX.X & XX.X\,/\,XX.X \\
\bottomrule
\end{tabular}
\end{table}

% ---------------------------------------------------------------------------
\subsection{Route Analysis}
\label{sec:exp:routes}
% ---------------------------------------------------------------------------

To understand \emph{how} the consensus router allocates compute, we
analyze the behavior of each difficulty route in detail.

\paragraph{Per-route accuracy.}
Table~\ref{tab:route-accuracy} breaks down accuracy by assigned route
for the best \method{} configuration ($K = 4$, $B_{\mathrm{probe}} = 128$).
Easy-route samples achieve the highest accuracy (XX.X\%), confirming
that the consensus signal reliably identifies problems the model can
solve with minimal compute.
Medium-route accuracy (XX.X\%) is intermediate, while hard-route accuracy
(XX.X\%) is lowest---as expected, since these are the problems where the
model genuinely struggles.
Critically, the easy-route accuracy exceeds what these same samples
achieve under Fixed-512 (XX.X\%), demonstrating that extended reasoning
\emph{hurts} on easy problems (the overthinking phenomenon discussed
in \S\ref{sec:intro}).

\begin{table}[t]
\centering
\caption{%
  \textbf{Per-route analysis} for \method{} ($K = 4$,
  $B_{\mathrm{probe}} = 128$, $n = 200$).
  ``Fixed-512 acc on same'' reports the accuracy of Fixed-512 on the
  subset of samples assigned to each route.
}
\label{tab:route-accuracy}
\small
\begin{tabular}{@{}lcccc@{}}
\toprule
\textbf{Route} & \textbf{\% Samples} & \textbf{Acc (\%)} &
  \textbf{Avg Tok} & \textbf{Fixed-512 acc on same (\%)} \\
\midrule
Easy   & XX.X & XX.X & XX.X & XX.X \\
Medium & XX.X & XX.X & XX.X & XX.X \\
Hard   & XX.X & XX.X & XX.X & XX.X \\
\midrule
Overall & 100.0 & XX.X & XX.X & XX.X \\
\bottomrule
\end{tabular}
\end{table}

\paragraph{Token savings per route.}
The easy route consumes only $K \times B_{\mathrm{probe}}$ tokens
(XX.X on average), saving XX.X\% relative to Fixed-512.
The medium route consumes an average of XX.X tokens (exploration +
extension), while the hard route consumes XX.X tokens (exploration +
full deliberation).
Across all routes, the weighted average token consumption is XX.X,
representing a XX.X\% reduction compared to Fixed-512.

\paragraph{Consensus score and correctness.}
Figure~\ref{fig:consensus-correct} plots the relationship between the
consensus confidence score $\gamma$ and per-sample correctness.
Samples with $\gamma \geq 0.8$ are correct XX.X\% of the time, while
samples with $\gamma < 0.4$ are correct only XX.X\% of the time.
The Spearman rank correlation between $\gamma$ and binary correctness
is $\rho = $ XX.XX ($p < $ XX), empirically validating
Proposition~\ref{prop:monotone}'s prediction that consensus is a
monotone difficulty estimator.
This strong correlation explains why the simple three-tier routing rule
in Eq.~\eqref{eq:routing} is effective: the consensus signal separates
difficulty levels well enough that even coarse binning into three routes
captures most of the available adaptive gain.

\begin{figure}[t]
\centering
% \includegraphics[width=0.55\textwidth]{figures/consensus_vs_correctness.pdf}
\fbox{\parbox{0.55\textwidth}{\centering\vspace{2cm}%
  \textit{Placeholder: scatter/binned plot of consensus score $\gamma$
  vs.\ correctness rate, with fitted logistic curve.}%
\vspace{2cm}}}
\caption{%
  \textbf{Consensus score vs.\ correctness.}
  Each point represents a bin of samples grouped by consensus confidence
  $\gamma$; the $y$-axis shows the fraction of correct answers.
  The strong positive correlation ($\rho =$ XX.XX) confirms that
  multi-path consensus is a reliable difficulty signal.
}
\label{fig:consensus-correct}
\end{figure}
